\subsection{Jet bundles and prolongation}

We recall in this subsection some well-known theory of jet bundles, drawn primarily from~\cite{varbi, saunders}. Although the reader is likely familiar with this theory already, we include the following outline both to introduce our rather unconventional notation (which we choose for consistency with the bundles of germs to be introduced in Section~\ref{pseudobundlesofgerms}) and to point out the structures that will be~most relevant in our constructions.

\begin{Definition}
Let $\pi_{B}\colon B\rightarrow M$ be a fibre bundle, with fibre $F$, and let $k\geq0$. We~say that two local sections $\sigma_{1}$ and $\sigma_{2}$ of $\pi_{B}$ defined in a neighbourhood of $x\in M$ have the same~$k$-\textit{jet at}~$x$ if $\sigma_{1}(x) =b= \sigma_{2}(x)$, and for any local coordinate neighbourhood $\OO\cong\OO_{M}\times\OO_{F}$ of $b$, one~has
\begin{gather*}
\frac{\partial^{|I|}\sigma^{i}_{1}}{\partial x^{I}}(x) = \frac{\partial^{|I|}\sigma^{i}_{2}}{\partial x^{I}}(x)
\end{gather*}
for all $i = 1,\dots,\dim(F)$, and all multi-indices $I$ with $|I|\leq k$. Having the same $k$-jet at a point~$x$ is an equivalence relation on the set of local sections defined about $x$, and we denote the $k$-jet equivalence class of any such local section by $[\sigma]^{k}_{x}$.
\end{Definition}

The $k$-jets of local sections fit into a fibre bundle in a natural way.

\begin{Definition}
Let $\pi_{B}\colon B\rightarrow M$ be a fibre bundle, and let $k\geq 0$. For each $x\in M$, denote the set of all $k$-jets of local sections defined near $x$ by $\Gamma_{k}(\pi_{B})_{x}$, and define
\begin{gather*}
\Gamma_{k}(\pi_{B}):=\bigsqcup_{x\in M}\Gamma_{k}(\pi_{B})_{x},
\end{gather*}
with $\pi_{B}^{k}\colon \Gamma_{k}(\pi_{B})\rightarrow M$ denoting the canonical projection. Define $\pi_{B}^{0,k}\colon \Gamma_{k}(\pi_{B})\rightarrow B$ by
\begin{gather*}
\pi_{B}^{0,k}\big(x,[\sigma]^{k}_{x}\big):=\sigma(x),
\end{gather*}
and observe that any choice of local coordinate trivialisation $(x^{i},f^{\alpha})$ on $\OO$ about a point $b\in B$ determines coordinates
\begin{gather*}%\label{coordsjets}
\big(x^{i},f^{\alpha},f^{\alpha}_{i},\dots,f^{\alpha}_{I}\big):=\bigg(x^{i},f^{\alpha},\frac{\partial f^{\alpha}}{\partial x^{i}},\dots,\frac{\partial^{|I|}f^{\alpha}}{\partial x^{I}}\bigg)
\end{gather*}
on the set $\big(\pi^{0,k}_{B}\big)^{-1}(\OO)$. These coordinates give the projection $\pi_{k}\colon \Gamma_{k}(\pi_{B})\rightarrow M$ the structure of~a~fibre bundle, called the $k^{\rm th}$ \textit{order jet bundle} of $\pi_{B}$.
\end{Definition}

The jet bundles of a fibre bundle $\pi_{B}\colon B\rightarrow M$ admit projections $\pi_{B}^{k,l}\colon \Gamma_{l}(\pi_{B})\rightarrow \Gamma_{k}(\pi_{B})$ defined by
\begin{gather*}
\pi_{B}^{k,l}\big([\sigma]^{l}_{x}\big):=[\sigma]^{k}_{x},\qquad [\sigma]^{l}_{x}\in \Gamma_{l}(\pi_{B})
\end{gather*}
for any $l\geq k$, and these projections form a projective system. The~projective limit of this system, denoted $\Gamma_{\infty}(\pi_{B})$, inherits a natural smooth structure as a projective limit of manifolds \cite[Chapter One, Section A]{varbi} (equivalently via the projective limit diffeology~\cite[Section 1.39]{diffeology}). The~space $\Gamma_{\infty}(\pi_{B})$ is usually identified with the set of $\infty$-jets of local sections of $\pi_{B}$, and admits a canonical projection $\pi_{B}^{\infty}\colon \Gamma_{\infty}(\pi_{B})\rightarrow M$ given by
\begin{gather*}
\pi_{B}^{\infty}\big([\sigma]^{\infty}_{x}\big):=x,\qquad [\sigma]^{\infty}_{x}\in \Gamma_{\infty}(\pi_{B}).
\end{gather*}
One therefore obtains a hierarchy of jet bundles
\begin{gather*}
\Gamma_{\infty}(\pi_{B})\rightarrow\cdots\xrightarrow{\pi_{B}^{k,k+1}} \Gamma_{k}(\pi_{B})\xrightarrow{\pi_{B}^{k-1,k}}\cdots\xrightarrow{\pi_{B}^{1,2}} \Gamma_{1}(\pi_{B})\xrightarrow{\pi_{B}^{0,1}}B\xrightarrow{\pi_{B}}M.
\end{gather*}

Since we will be concerned primarily with singular foliations, which arise from families of~vec\-tor fields, we will need to know about vector fields on jet bundles. A~particularly important class of vector fields on fibre bundles, in which we will be primarily interested, is those that are \textit{projectable} in the following sense.

\begin{Definition}%\label{projectable}
Let $\pi_{B}\colon B\rightarrow M$ be a fibre bundle. A~vector field $X$ on $B$ is said to be \textit{pro\-jec\-table} if there is a vector field $(\pi_{B})_{*}(X)$ on $M$ for which
\begin{gather}\label{projectableeqn}
{\rm d}\pi_{B}\circ X = (\pi_{B})_{*}(X)\circ \pi_{B}
\end{gather}
on all of $B$. We~denote by $\XF_{\proj}(B)$ the set of projectable vector fields on $B$.
\end{Definition}

Projectable vector fields on a bundle prolong in a natural way to vector fields on the associated jet bundles.

\begin{Definition}
Let $\pi_{B}\colon B\rightarrow M$ be a fibre bundle, and let $X\in\XF_{\proj}(B)$ be a (possibly time-dependent) projectable vector field. For $1\leq k<\infty$, the \textit{$k$-jet prolongation} of $X$ is the vector field $\pf^{k}(X)\in\XF(\Gamma_{k}(\pi_{B}))$ defined by
\begin{gather*}
\pf^{k}(X)\big(x,[\sigma]^{k}_{x}\big):=\frac{\rm d}{{\rm d}t}\bigg|_{t=0}\Big(\Fl^{(\pi_{B})_{*}X}_{t,0}(x),\big[\Fl^{X}_{t,0}\circ \sigma\circ\Fl^{(\pi_{B})_{*}X}_{0,t}\big]^{k}_{\Fl^{(\pi_{B})_{*}X}_{t,0}(x)}\Big)
\end{gather*}
for all $[\sigma]^{k}_{x}\in \Gamma_{k}(\pi_{B})$. We~denote by $\XF_{\proj}(\Gamma_{k}(\pi_{B}))$ the image of $\pf^{k}$.
\end{Definition}

It will be useful later to note here that if $X$ is a time-independent projectable vector field on~a~bundle $\pi_{B}$, it is given in coordinates by $X = a^{i}\frac{\partial}{\partial x^{i}} + b^{\alpha}\frac{\partial}{\partial f^{\alpha}}$, where $a^{i}$ are smooth functions depending only on the $x^{i}$ while the $b^{\alpha}$ depend on both the $x^{i}$ and the $f^{\alpha}$. Then for any $1\leq k<\infty$, the $k$-jet prolongation $\pf^{k}(X)$ of $X$ is given in coordinates over a point $(x^{i},f^{\alpha},f^{\alpha}_{i},\dots)\in \Gamma_{k}(\pi_{B})$ by the formula (cf.~\cite[Theorem~2.36]{olver})
\begin{gather}\label{prolongation}
\pf^{k}(X) = a^{i}D_{i}^{(k)}+\sum_{|I|=k}\!\bigg(D^{(k)}_{I}b^{\alpha}-\sum_{J\subset I}\big(D^{(k)}_{I\setminus J}a^{i}\big)f^{\alpha}_{Ji}\bigg)\frac{\partial}{\partial f^{\alpha}_{I}}+\!\sum_{|I|=0}^{k-1}D^{(k)}_{I}(b^{\alpha}-a^{i}f^{\alpha}_{i})\frac{\partial}{\partial f^{\alpha}_{I}},
\end{gather}
where the sum over $J\subset I$ is a sum over all \textit{strict} subsets of the multi-index $J$, and where we have absorbed the constants arising from the symmetry of mixed partial derivatives into our notation as in~\cite[equation~(1.15)]{varbi}. Here $D^{(k)}_{I} = D^{(k)}_{i_{1}}\cdots D^{(k)}_{i_{k}}$, where $D^{(k)}_{i}$ is the \textit{total derivative operator} defined by $D^{(k)}_{i} = \frac{\partial}{\partial x^{i}}+\sum_{|I|=0}^{k-1}f^{\alpha}_{Ii}\frac{\partial}{\partial f^{\alpha}_{I}}$.

The argument of Proposition~\ref{germprol} can be used to deduce the following fact, which justifies our choice of notation in denoting the image of $\pf^{k}$ in $\XF(\Gamma_{k}(\pi_{B}))$ by $\XF_{\proj}(\Gamma_{k}(\pi_{B}))$.

\begin{Proposition}
Let $\pi_{B}\colon B\rightarrow M$ be a fibre bundle, and let $X$ be a projectable vector field on~$B$. Then the $\pf^{k}(X)$ are a projectable family, in the sense that
\begin{gather*}
{\rm d}\pi_{B}^{l,k}\circ \pf^{k}(X) = \pf^{l}(X)\circ\pi_{B}^{l,k}
\end{gather*}
on $\Gamma_{k}(\pi_{B})$ for all $l\leq k$.
\end{Proposition}

Projectability of the $\pf^{k}(X)$ guarantees (cf.~\cite[equation~(1.11)]{varbi}) that pointwise they admit a~pro\-jective limit, defining a section $\pf^{\infty}(X)$ of the projective limit tangent bundle of $\Gamma_{\infty}(\pi_{B})$ (cf.~\cite[Chapter One, Section B]{varbi}) for which the following holds.

\begin{Proposition}\label{towerfield1}
Let $\pi_{B}\colon B\rightarrow M$ be a fibre bundle. Then for each $l\leq k\leq\infty$ there is a~bracket-preserving homomorphism $\big(\pi_{B}^{l,k}\big)_{*}\colon \XF_{\proj}(\Gamma_{k}(\pi_{B}))\rightarrow\XF_{\proj}(\Gamma_{l}(\pi_{B}))$ and the diagram
\begin{center}
\begin{tikzcd}[column sep = huge]
& \XF_{\proj}(\Gamma_{\infty}(\pi_{B})) \ar[d]
\\
& \vdots \ar[d,"\left(\pi_{B}^{1,2}\right)_{*}"]
\\
& \XF_{\proj}(\Gamma_{1}(\pi_{B})) \ar[d,"\left(\pi_{B}^{0,1}\right)_{*}"]
\\
\XF_{\proj}(B) \ar[uuur,"\pf^{\infty}"] \ar[ur,"\pf^{1}"] \ar[r,"\id"] & \XF_{\proj}(B)
\end{tikzcd}
\end{center}
commutes.
\end{Proposition}

\subsection{Diffeology}

We recall in this subsection some basic objects of study in diffeology that will be relevant for our constructions. The~most comprehensive reference on diffeology is the wonderful book~\cite{diffeology} by P. Iglesias-Zemmour.

\begin{Definition}\label{diffeology}
A function $\varphi\colon U\rightarrow \mathcal{X}$ from an open subset $U$ of some finite-dimensional Euclidean space to a set $\XX$ is called a \textit{parametrisation}. A~\textit{diffeology} on a set $\XX$ is a family D of~para\-metrisations satisfying the following axioms.
\begin{enumerate}\itemsep=0pt
\item The family D contains all constant parametrisations.
\item If $\varphi\colon U\rightarrow \XX$ is a parametrisation such that every point $u\in U$ has an open neighbourhood $V\subset U$ for which $\varphi|_{V}$ is an element of D, then $\varphi$ itself is an element of D.
\item For every element $\varphi\colon U\rightarrow \XX$ of D, every open set $V$ of any finite-dimensional Euclidean space, and for every smooth function $f\colon V\rightarrow U$, the composite $\varphi\circ f\colon V\rightarrow \XX$ is contained in D.
\end{enumerate}
A set with a diffeology is called a \textit{diffeological space}, and the elements of the diffeology are called its \textit{plots}. If $\XX$ and $\YY$ are two diffeological spaces, then a function $f\colon \XX\rightarrow \YY$ is said to be \textit{smooth} if for every plot $\varphi\colon U\rightarrow \XX$ of $\XX$, the composite $f\circ\varphi\colon U\rightarrow \YY$ is a plot of $\YY$. A~smooth bijection of diffeological spaces is said to be a \textit{diffeomorphism} if it is smooth with smooth inverse.
\end{Definition}

Every manifold is a diffeological space, with diffeology constituted by the set of all parametrisations that are smooth in the usual sense. Moreover a map between manifolds is smooth in the manifold sense if and only if it is smooth in the diffeological sense. Thus the category of~mani\-folds and smooth maps is a full and faithful subcategory of the category of diffeological spaces and smooth maps.

Diffeologies can be pushed forward and pulled back by functions of sets. This fact will be invoked frequently for our constructions.

\begin{Definition}
Let $\XX$ and $\YY$ be sets, and let $f\colon \XX\rightarrow \YY$ be a function.
\begin{enumerate}\itemsep=0pt
\item If $\XX$ has a diffeology, then the \textit{pushforward diffeology induced by $f$} is defined by declaring a parametrisation $\varphi\colon U\rightarrow \YY$ to be a plot if and only if every $u\in U$ has an open neighbourhood $V\subset U$ such that either $\varphi|_{V}$ is constant, or equal to the composite $f\circ\psi$ for some plot $\psi\colon V\rightarrow \XX$ of $\XX$. The~map $f$ is said to be a \textit{subduction} if it is surjective and if $\YY$ is equipped with the pushforward diffeology induced by $f$.
\item If $\YY$ has a diffeology, then the \textit{pullback diffeology induced by $f$} is defined by declaring a~parametrisation $\psi\colon U\rightarrow \XX$ to be plot if and only if the composite $f\circ\psi\colon U\rightarrow \YY$ is a~plot of~$\YY$. The~map $f$ is said to be an \textit{induction} if it is injective and if $\XX$ is equipped with the pullback diffeology from $\YY$.
\end{enumerate}
The following special cases are of particular importance. Let $\XX$ be a diffeological space.
\begin{enumerate}\itemsep=0pt
\item If $\sim$ is any equivalence relation on $\XX$, then the \textit{quotient diffeology} $\XX/{\sim}$ is the pushfoward diffeology arising from the quotient $\XX\rightarrow \XX/{\sim}$.
\item If $S$ is any subset of $\XX$, then the \textit{subspace diffeology} on $S$ is the pullback diffeology arising from the inclusion $S\hookrightarrow \XX$.
\item If $\YY$ is any other diffeological space, then the \textit{product diffeology} on $\XX\times \YY$ is the smallest diffeology for which the projections onto the factors are subductions.
\end{enumerate}
Quotients, subspaces and products will always be assumed to be equipped with the respective diffeologies defined above unless otherwise stated.
\end{Definition}

One of the features of the category of diffeological spaces is that the set of all morphisms between any two objects in the category is itself an object.

\begin{Definition}
Let $\XX$ and $\YY$ be diffeological spaces, and denote by $C^{\infty}(\XX,\YY)$ the set of~all smooth maps $f\colon \XX\rightarrow \YY$. The~\textit{functional diffeology} on $C^{\infty}(\XX,\YY)$ is defined by declaring a~para\-metrisation $\tilde{f}\colon U\rightarrow C^{\infty}(\XX,\YY)$ to be a plot if and only if the associated map $U\times\XX\ni(u,x)\mapsto\tilde{f}(u)(x)\rightarrow\YY$ is smooth.
\end{Definition}

The familiar notion of a fibre bundle over a manifold has a far reaching generalisation to diffeological spaces. It is the flexibility afforded by this generalisation that permits most of the constructions in this paper.

\begin{Definition}%\label{pb}
A \textit{diffeological pseudo-bundle} is a subduction $\pi_{\BB}\colon \BB\rightarrow \XX$ of diffeological spaces. Such a pseudo-bundle is in particular called a \textit{diffeological vector pseudo-bundle} if each fibre of $\pi_{\BB}$ is a vector space for which the vector space operations are smooth with respect to the subspace diffeology, and such that the zero section is smooth.
\end{Definition}

A diffeological pseudo-bundle need not have fibres that are all diffeomorphic, and of course need not be locally trivial in any sense (see for instance Example~\ref{germfol}). An important subclass of~diffeological pseudo-bundles are \textit{diffeological bundles}, which have mutually diffeomorphic fibres and which are locally trivial under pullbacks by plots. Diffeological bundles are distinguished by the behaviour of their \textit{structure groupoids}. Before we give the definition of this object, let us record what we mean by diffeological categories and groupoids.

\begin{Definition}
Let $\CC$ be a small category, with object set identified as a subset of the morphism set via the map which sends each object to its associated identity morphism. We~say that~$\CC$ is a \textit{diffeological category} if its set of morphisms is equipped with a diffeology for which the range, source, and composition are all smooth. If $\CC$ is in addition a groupoid, whose inversion map is smooth, we call $\CC$ a \textit{diffeological groupoid}.
\end{Definition}

Let now $\pi_{\BB}\colon \BB\rightarrow \XX$ be a smooth surjection of diffeological spaces. Denote by $\Aut(\pi_{\BB})$ the groupoid with object set $\XX$, and with morphisms from $x$ to $y$ constituted by the set $\Diff(\BB_{x},\BB_{y})$ of all diffeomorphisms from the fibre $\BB_{x}$ over $x$ to the fibre $\BB_{y}$ over $y$, with the obvious range, source, inversion and composition. The~groupoid $\Aut(\pi_{\BB})$ admits a smallest diffeology, called the \textit{functional diffeology}, under which the evaluation map $\ev\colon \Aut(\pi_{\BB})\times_{s,\pi_{\BB}}\BB\ni(f,b)\mapsto f(b)\in \BB$ is smooth, and under which $\Aut(\pi_{\BB})$ is a diffeological groupoid (see~\cite[Section~8.7]{diffeology} for details).

\begin{Definition}\label{structuregpd}
 Let $\pi_{\BB}\colon \BB\rightarrow \XX$ be a surjection of diffeological spaces. Equipped with the functional diffeology, we refer to $\Aut(\pi_{\BB})$ as the \textit{structure groupoid} of the surjection $\pi_{\BB}$. We~say that $\pi_{\BB}$ is a \textit{diffeological bundle} if the characteristic map $(r,s)\colon \Aut(\pi_{\BB})\rightarrow \XX\times \XX$ is a sub\-duction.
\end{Definition}

Diffeological bundles, unlike general diffeological pseudo-bundles, have a typical fibre to which all other fibres are diffeomorphic, and the pullback of a diffeological bundle along any plot is locally trivial~\cite[p.~240]{diffeology}.

By definition, singular foliations arise from certain families of sections of tangent bundles. To~use diffeology to study singular foliations, therefore, we need a notion of tangent bundle for a diffeological space. A~number of definitions have been proposed for this purpose, which, while coincident for manifolds, do not coincide for general diffeological spaces (see~\cite{cw} for a detailed discussion). The~point of view that we find useful here, as in~\cite{mac3}, is that of \textit{internal} tangent spaces and bundles. The~paper~\cite{cw} provides a categorical definition of internal tangent spaces based on work of Hector~\cite{hector}, which may be summarised as follows.

\begin{Definition}\label{tangentbundle}
Let $\XX$ be a diffeological space and let $x\in \XX$. Denote by $p_{x}$ the set of all \textit{plots centered at $x$}, that is, plots $\varphi\colon U\rightarrow \XX$ such that $0\in U$ and $\varphi(0) = x$. Denote by $T_{0}\dom(\varphi)$ the tangent space at zero of the domain of any such plot, and by $\vec{v}_{\varphi}$ the image of $\vec{v}\in T_{0}\dom(\varphi)$ in the direct sum $\bigoplus_{\varphi\in p_{x}}T_{0}\dom(\varphi)$. The~\textit{internal tangent space of $\XX$ at $x$} is the quotient space~$T_{x}\XX$ of the direct sum
\begin{gather*}
\bigoplus_{\varphi\in p_{x}} T_{0}\dom(\varphi)
\end{gather*}
by the subspace generated by all vectors of the form $f_{*}(\vec{v})_{\varphi} - \vec{v}_{\varphi\circ f}$, where $\varphi\in p_{x}$, $f\colon \dom(f)\rightarrow\dom(\varphi)$ is any smooth map of open Euclidean domains with $0\in\dom(f)$ and $f(0) = 0$, and with $\vec{v}\in T_{0}\dom(f)$. The~class of an element $\vec{v}_{\varphi}$, $\vec{v}\in T_{0}\dom(\varphi)$, will be denoted $\varphi_{*}(\vec{v})$. By~\cite[Proposition~3.3]{cw}, elements of $T_{x}\XX$ can always be written as linear combinations of $\varphi_{*}({\rm d}/{\rm d}t)$ for 1-plots $\varphi\colon (-\epsilon,\epsilon)\rightarrow\XX$, which we will usually write as ${\rm d}/{\rm d}t|_{0}\varphi_{t}$.
\end{Definition}

Let $\XX$ be a diffeological space, and consider the set $T\XX:=\bigsqcup_{x\in \XX}T_{x}\XX$. For any plot \mbox{$\varphi\colon U\rightarrow \XX$} and for any $u\in U$, denote by $\tau_{u}\colon v\mapsto v+u$ the $u$-translation map, so that $\tau_{u}^{-1}(U)$ contains $0$ and $\varphi\circ\tau_{u}\colon \tau_{u}^{-1}(U)\rightarrow \XX$ is a plot centered at $\varphi(u)$. Define ${\rm d}\varphi\colon TU\rightarrow T\XX$ by the formula
\begin{gather*}
{\rm d}\varphi(u,\vec{v}):=(\varphi(u),(\varphi\circ\tau_{u})_{*}(\vec{v})),\qquad (u,\vec{v})\in TU.
\end{gather*}
These maps were first considered by Hector~\cite{hector}. Then there exists a smallest diffeology on $T\XX$, called the \textit{dvs diffeology}~\cite{cw}, for which the natural projection $\pi_{T\XX}\colon T\XX\rightarrow \XX$ is a diffeological vector pseudo-bundle, and which contains the parametrisations ${\rm d}\varphi\colon TU\rightarrow T\XX$ as plots.

\begin{Definition}
Let $\XX$ be a diffeological space. The~diffeological vector pseudo-bundle $\pi_{T\XX}$: $T\XX\rightarrow \XX$ is called the \textit{internal tangent bundle of $\XX$}.
\end{Definition}

The internal tangent bundle is functorial under smooth maps of diffeological spaces.

\begin{Definition}\label{differential}
Let $f\colon \XX\rightarrow \YY$ be a smooth map of diffeological spaces. For any plot $\varphi$ of $\XX$ centered at $x$, define
\begin{gather*}
{\rm d}f(x,\varphi_{*}(\vec{v})):=(f(x),(f\circ\varphi)_{*}(\vec{v})),\qquad x\in \XX
\end{gather*}
and extend by linearity to a map ${\rm d}f\colon T\XX\rightarrow T\YY$. Then ${\rm d}f$ is a smooth map~\cite[Proposition~4.8]{cw} called the \textit{pushforward} or \textit{differential} of $f$.
\end{Definition}
